\scanpage{123}
\mspage{114}
% archive sheet number: 123

\uline{Proposition} Une préfigure quelconque est déterminée par les données suivantes

\unsure{4°|1°}) Une figure \textcircled{$\Psi$}\inserted{sans sommets isolés et} à composantes connexes des \deleted{multi} segments, \inserted{\deleted{(i.e. $\delta\Psi=\partial\Phi$} i.e. telle que $\delta\Psi=\partial\Phi$)} déterminée par
\[
D=(X_1\lhd X_2\lhd\deleted{\unsure{X_3}}\ldots X_n)\in\mathrm{Drap}^{*}(\unsure{\mathcal{M}|M}_1).
\]
\deleted{2°) Une}
\deleted{tel} que
\[
\Psi=\Phi_D=\bigcup_{X\in D}\widetilde{X}_{\deleted{\unsure{i}}}=\bigcup_{1\leq i\leq n}\widetilde{X}_i
\]
(réunion \unsure{disjoint|disjointe}). On aura
\[
\partial\Psi=\delta\Psi=\bigcup_{X\in D}\partial X=\bigcup_{1\leq i\leq n}\partial X_i
\]
(réunion disjointe), soit
\[
\delta\Psi=\{a_1,b_1,a_2,b_2,\ldots,a_n,b_n\}
\]
si
\[
\partial X_i=\{a_i,b_i\}\,.
\]

On peut dire aussi que $\Psi$ est déterminée par une suite \unsure{strict\textsuperscript{t}} croissante de cardinal pair
\[
a_1<b_1<a_2<b_2\ldots<a_n<b_n\,,
\]
et on aura
\[
X_i=\unsure{S|\mathscr{S}}_{a_i,b_i}\qquad 1\leq i\leq n.
\]
% slanted
\margin{2°) ici, \unsure{cf. p°|cf. V°} \unsure{suivante|suivant}\par\tikz{\draw (0,0)--(1.8,0);}}

3°) Un drapeau \textcircled{$\Delta$}$\in\mathrm{Drap}^{*}(\unsure{\mathcal{L}|L})$
\[
\Delta\supset\partial\Psi
\]
\deleted{3°) Une partie $\Delta_{\mathrm{lac}}$}

d'où on déduit une partie
